% Auto-generated by notebooks/99_markov_vs_age_dependent.py -- do not edit by hand.
\begin{table}[htbp]
\centering
\begin{tabular}{llrrrrr}
\toprule
Condition & Model & $k$ & $\log L$ & AIC & Vuong $z$ & Vuong $p$ \\
\midrule
DMSO\_day\_10 & Langevin (this work) & 3 & -404.6 & 815.2 & -- & -- \\
DMSO\_day\_10 & Gompertz & 2 & -342.9 & 689.8 & -7.61 & 2.6e-14 \\
Auxin\_day\_21 & Langevin (this work) & 0 & -459.6 & 919.3 & -- & -- \\
Auxin\_day\_21 & Gompertz & 0 & -442.0 & 884.0 & -2.04 & 0.041 \\
\bottomrule
\end{tabular}
\caption{Grouped-data maximum-likelihood fit of our Langevin model (parameters from \texttt{0\_auto\_fit\_parameters\_mle.py}, not refit here) versus a matched-complexity Gompertz hazard, per condition. For Auxin\_day\_21, both models' phase 1 and phase 2 parameters are fixed from other conditions' own fits (DMSO\_day\_10 for phase 1, Auxin\_day\_10 for phase 2), so $k=0$ parameters are estimated from Auxin\_day\_21's own data in either model -- a genuine out-of-sample prediction, not a fit. $\Delta\mathrm{AIC} = \mathrm{AIC}_\mathrm{Gompertz} - \mathrm{AIC}_\mathrm{Langevin} = -125.4$ (DMSO\_day\_10), $-35.3$ (Auxin\_day\_21); negative values favor Gompertz. Vuong $z$/$p$ (reported on the Gompertz row of each condition, testing that condition's Langevin-vs-Gompertz pair) is Vuong's (1989) closeness test, AIC-corrected for $k$; $z>0$ favors Langevin. For DMSO\_day\_10, Gompertz is significantly closer to the true generating process ($p=2.6e-14$). For Auxin\_day\_21, Gompertz is significantly closer to the true generating process ($p=0.041$).}
\label{tab:supp_note2_model_comparison}
\end{table}
